\section{Experimental Study}
\label{sec:experiments}

\subsection{Scenario tree}
\label{sec:scentree}
We use a two-stage design: one-factor-at-a-time (OFAT) screening to rank the
levers, followed by a combined run on the two highest-impact levers, then a
demand sweep for 2030 and the Industry~4.0 scrap scenario.

\begin{figure}[H]
\centering
\begin{forest}
  for tree={
    grow'=0, draw, rounded corners, font=\footnotesize,
    edge={-Stealth}, l sep=10mm, s sep=2mm, anchor=west,
    child anchor=west, parent anchor=east, align=left}
  [Baseline\\(dryer = bottleneck)
    [OFAT screening
      [{+1 Dryer ($n_\text{dryers}=2$)}]
      [{+1 Cleaning team ($c=2$)}]
      [{+1 Packaging line ($n_\text{pack}=4$)}]
      [Resequence (group by product)]
    ]
    [Best combination\\(dryer$+$clean team)
      [{dryer$=2$, clean$=2$}]
      [{dryer$=2$, clean$=2$, pack$=4$}]
    ]
    [2030 demand sweep\\($\times$1.25\,/\,$\times$1.5\,/\,$\times$2.0)
      [Max sustainable throughput]
    ]
    [Industry 4.0\\(scrap $\rightarrow$ 4\,\%)]
  ]
\end{forest}
\caption{Scenario tree for the experimental study.}
\label{fig:scentree}
\end{figure}

\subsection{OFAT results and bottleneck shift}
\label{sec:btshift}
Table~\ref{tab:btshift} reports mean station utilisation across all scenarios.
The bold entry in each row is the binding constraint (highest utilisation).
The key finding is that adding a dryer alone shifts the bottleneck immediately
to the cleaning team, revealing that the two resources are \emph{co-constraints}:
neither can be relieved without also addressing the other.

\begin{table}[H]
\centering
\caption{Mean station utilisation (\%) by scenario ($N=20$ reps each).
         Bold = binding constraint. Throughput in bottles shipped (mean\,$\pm$\,95\,\% CI).}
\label{tab:btshift}
\small
\begin{tabularx}{\textwidth}{l c c c c r}
\toprule
Scenario & Dryer & Clean team & Packaging & Press &
  Throughput (bottles) \\
\midrule
Baseline
  & \textbf{94.6} & 81.9 & 54.5 & 6.5
  & $67\,487 \pm 1\,275$ \\
+1 dryer
  & 76.3 & \textbf{95.2} & 74.0 & 6.6
  & $67\,701 \pm 1\,074$ \\
+1 cleaning team
  & 90.7 & \textbf{91.2} & 88.1 & 6.6
  & $67\,894 \pm 1\,151$ \\
+1 packaging line
  & \textbf{94.8} & 82.1 & 41.0 & 6.6
  & $67\,919 \pm 1\,395$ \\
Resequence (grouped)
  & \textbf{93.2} & 86.5 & 54.9 & 6.5
  & $67\,579 \pm 1\,387$ \\
\midrule
+1 dryer \& +1 clean team
  & 58.3 & \textbf{94.2} & 93.6 & 6.6
  & $68\,794 \pm 904$ \\
+1 dryer, clean \& pack line
  & 72.9 & \textbf{93.5} & 77.9 & 6.6
  & $67\,633 \pm 1\,320$ \\
\bottomrule
\end{tabularx}
\end{table}

No individual OFAT change produces a statistically significant throughput gain:
all single-lever confidence intervals overlap with the baseline. Only the
combined dryer\,+\,clean-team scenario approaches significance
($\Delta = +1{,}307$ bottles; CI $[67{,}890$--$69{,}698]$ vs.\ baseline
$[66{,}212$--$68{,}762]$, marginal overlap of 872 bottles).
Adding a fourth packaging line on top of the combined scenario \emph{reduces}
throughput back to baseline levels: the extra line generates additional
cleaning-team demand (setup and changeover), overloading the already-stressed
cleaning team.

\subsection{Results: throughput across scenarios}
\label{sec:expresults}
% TODO: insert bar chart with 95 % CI error bars (one bar per scenario).
\begin{figure}[H]
\centering
\fbox{\parbox[c][40mm][c]{0.7\textwidth}{\centering\itshape
Throughput by scenario (with 95\,\% CI error bars) --- placeholder.}}
\caption{Headline KPI comparison across scenarios.}
\label{fig:expresults}
\end{figure}

\subsection{2030 demand and Industry 4.0}
\label{sec:demand40}
Using the calibrated baseline and the best configuration (dryer$=2$,
clean-team$=2$), we scale demand via \texttt{demandMultiplier} to locate the
maximum weekly volume the line can clear before backlog grows unbounded.

\begin{table}[H]
\centering
\caption{Mean bottles shipped under rising demand ($N=20$ reps).
         \emph{TODO: fill after demand-sweep runs.}}
\label{tab:demand}
\begin{tabularx}{\textwidth}{l c c c}
\toprule
Configuration & $\times$1.25 & $\times$1.5 & $\times$2.0 \\
\midrule
Baseline         & [\,] & [\,] & [\,] \\
Best combination & [\,] & [\,] & [\,] \\
\bottomrule
\end{tabularx}
\end{table}

\paragraph{Industry 4.0 --- predictive maintenance.}
Reducing premix scrap from the data-derived \(\approx\)6\,\% to 4\,\%
via \texttt{premixFailureProbability\,=\,0.04} models the effect of
sensors that detect blender anomalies before a batch is lost.
Because fewer scrapped batches means more tablets reaching the dryer and
packaging constraint, the net throughput effect depends on whether the
constraint can absorb the extra load.
\emph{TODO: fill after Industry~4.0 run.} Result: [\,].

\subsection{Managerial value: benefit vs.\ effort}
\label{sec:value}
\begin{table}[H]
\centering
\caption{Each lever's mean throughput gain (vs.\ baseline 67\,487 bottles)
         against implementation effort. \emph{None of the single-lever gains
         are statistically significant; the combined gain approaches
         significance.}}
\label{tab:value}
\begin{tabularx}{\textwidth}{l c c X}
\toprule
Lever & $\Delta$ bottles & $\Delta$\,\% & Implementation effort / cost \\
\midrule
Resequence (group by product)
  & \pos{92}   & \pos{0.1\,\%}
  & Free --- scheduling change only \\
$+1$ dryer
  & \pos{214}  & \pos{0.3\,\%}
  & Medium --- equipment purchase, moderate lead time \\
$+1$ cleaning team member
  & \pos{407}  & \pos{0.6\,\%}
  & Low --- additional hire or shift \\
$+1$ packaging line
  & \pos{432}  & \pos{0.6\,\%}
  & High --- capex, long lead time (deferred to 2030 plan) \\
\midrule
\textbf{+1 dryer \& +1 clean team (combined)}
  & \pos{1{,}307} & \pos{1.9\,\%}
  & Medium --- only lever pair that relieves both co-constraints \\
\bottomrule
\end{tabularx}
\end{table}
